\documentclass[11pt]{article}

\usepackage[T1]{fontenc}
\usepackage[utf8]{inputenc}
\usepackage[english]{babel}
\usepackage{microtype}
\usepackage[margin=1in]{geometry}
\usepackage{xcolor}
\usepackage[hidelinks]{hyperref}

\definecolor{responseblue}{RGB}{0,70,140}

\title{Response to Reviewers\\[4pt]\large Isochrone overturning in steady glacier flow}
\author{Bradley Paul Lipovsky}
\date{}

\begin{document}
\maketitle

\section*{Reviewer 1}

\subsection*{R1C1: Applications to actual glaciers and folds}

\noindent\textbf{Comment.}
\begin{quote}
Although this work is glaciologically focused, I am not sure that there are sufficient applications to actual glaciers and their folds. The work is framed in terms of the Allan Hills, but for the unindoctrinated reader (or, in this case, referee), this context could be deepened. A schematic of the Allan Hills setting would be a useful start. Then returning to a real glacier after the Stokes numerical solution. That is, a calculation and larger contextualization than what is currently in the Discussion section. Are there radargrams that show the effects that the author is modeling? Including these types of connections to glacier system would ground the paper.
\end{quote}

\noindent\textcolor{responseblue}{\textbf{Response.} Thank you. We added a new Allan Hills context figure near the beginning of the Introduction. The first panel maps ALHIC1901 and ALHIC1902 on a 10~m Sentinel-2 image and locates Allan Hills within Antarctica. The other panels show the published \(^{40}\mathrm{Ar}_{\mathrm{atm}}\) age--depth relations for the two cores, with intervals of decreasing age with increasing depth highlighted. We also added text that describes the 1.5~km separation of the cores, their contrasting modelled flow regimes, and the continuity assumption needed to interpret an age reversal as an overturned continuous isochrone.}

\medskip

\noindent\textcolor{responseblue}{The revised Discussion returns to Allan Hills after the numerical example and relates the diagnostic to the near-stagnant basal flow, three-dimensional diversion around the nunatak, and multiple possible source regions inferred there. We did not add a radargram because the presently cited radar observations constrain spatial continuity and transport but do not provide a dated reflector that can be connected directly to the age reversals in these two cores. Instead, the revised figure uses the dated cores as the direct observational connection and states the continuity limitation explicitly.}

\medskip

\noindent\textcolor{responseblue}{We also did not add an Allan Hills-specific two-dimensional calculation. Existing reconstruction indicates strongly three-dimensional flow, so a two-dimensional site calculation could give a misleading test. The revised text states that evaluating steady overturning at Allan Hills requires a three-dimensional, multi-source material map.}

\subsection*{R1C2: Method of characteristics}

\noindent\textbf{Comment.}
\begin{quote}
It would be interesting to see an interpretation of the results that leans on the method of characteristics, e.g., along the lines of equations (57) and (58). It is operating in the background through the analysis of the Jacobian. I think that highlighting the connection could facilitate the discussion of the applicability of the results to a specific glacier system.
\end{quote}

\noindent\textcolor{responseblue}{\textbf{Response.} Thank you. We agree that the method of characteristics was already present in the analysis but was not identified clearly. We added text after the steady age equation to state that particle trajectories are its characteristics, age is elapsed travel time along them, and the deposition coordinate \(x_0\) labels them. In this interpretation, the order Jacobian is the sensitivity of the equal-age characteristic map to its initial surface label. No new governing equations are required because the existing Jacobian analysis already gives this characteristic sensitivity.}

\medskip

\noindent\textcolor{responseblue}{We also revised the trajectory-diagnostic section to identify Equations~(57) and (58) explicitly as a method-of-characteristics implementation. Equation~(57) traces a characteristic from the deposition boundary, Equation~(58) accumulates the signed shear history that controls separation from neighboring surface origins, and Equation~(59) gives the resulting Jacobian without finite differencing neighboring trajectories. The revised Discussion then states how this construction would be applied to a glacier: launch characteristics from candidate accumulation sources, follow them through a reconstructed velocity field to the core or radar transect, and evolve the source-patch Jacobian. At Allan Hills this application must be three-dimensional and allow multiple source regions.}

\end{document}
